\section{Congestion model}
\label{sec:congestion}

\begin{lede}
How crowded a cell is about to be, expressed as one number between 0 and 1 so
that it can modulate the score without ever dominating it.
\end{lede}

\subsection{Occupancy index}
\label{sec:congestion-index}

\code{OccupancyIndex} is rebuilt every timestep in $O(n)$ for $n$ robots: a
vertex-to-robot hash map, an aisle-to-count counter, and a spatial hash with
cell size $\max(1, R_{\mathrm{local}})$ for fast local-density queries. Aisle
load is a ratio, and may exceed 1:
\begin{equation}
  \mathrm{aisle\_load}(a) = \frac{\mathrm{occupancy}(a)}{\max(1,\ \mathrm{capacity}(a))}.
\end{equation}

\code{local\_density(v, exclude)} approximates
$|\{j : d_1(v, \pos{j}{t}) \leq R_{\mathrm{local}}\}|$, excluding one robot
(typically the querying one), in two passes: sum the $3 \times 3$ bucket window
around $v$ for a cheap upper bound, and if it is nonzero, refine exactly by
scanning the Manhattan diamond of radius $R_{\mathrm{local}} = 3$.

\subsection{Downstream lookahead}
\label{sec:congestion-downstream}

\begin{equation}
  \mathrm{downstream}(v, w)
    = \frac{1}{|\mathrm{tail}|}\sum_{u \in \mathrm{tail}} \ind{\text{$u$ occupied}},
  \qquad
  \mathrm{tail} = \mathrm{shortest\_route}(v, w, K)[1{:}],
\end{equation}
with $K = \param{downstream\_horizon} = 5$: the fraction of the next $K$ route
vertices beyond $v$ currently holding a robot. It is cached per $(v, w)$ per
timestep, since many robots query the same downstream segment.

\subsection{Congestion mixture}
\label{sec:congestion-mixture}

\begin{equation}
  C_i(v) = \frac{\omega_{\mathrm{loc}}\, \mathrm{local\_ratio}(v)
    + \omega_{\mathrm{aisle}}\, \mathrm{aisle\_load}(a(v))
    + \omega_{\mathrm{down}}\, \mathrm{downstream}(v, w_i)}
    {\omega_{\mathrm{loc}} + \omega_{\mathrm{aisle}} + \omega_{\mathrm{down}}}.
  \label{eq:congestion}
\end{equation}
All three weights default to $1.0$, so $C_i(v)$ is the mean of three ratios and
lies in $[0, 1]$. It returns $0$ unconditionally when
\param{congestion\_scoring} is off.

The first term is a \emph{ratio}, not a count:
\begin{equation}
  \mathrm{local\_ratio}(v)
    = \frac{\mathrm{local\_density}(v)}{|\{u \in V : d_1(u, v) \leq R_{\mathrm{local}}\}|},
\end{equation}
the fraction of reachable cells within the radius that hold a robot. The other
two were already ratios, so mixing a count with them made the local term
effectively the whole mixture and let $\mu C$ outrank the terms it is meant to
modulate (\Cref{sec:score-full}). \param{congestion\_normalisation}
\code{=False} restores the raw-count form.

\begin{caveat}
Normalising is a \emph{calibration} fix, not a performance one: on the bundled
maps its measured effect on throughput is neutral to slightly negative. What it
buys is that the weight ordering \Cref{eq:ordering} claims is the ordering the
system actually has at runtime. That is worth stating plainly rather than
presenting as a speed-up.
\end{caveat}

% worked-example: congestion
\begin{workedexample}
\textbf{Worked example.} The same robot and timestep. \code{C\_i(v)} mixes three occupancy signals with weights \code{ω\_local} = 1, \code{ω\_aisle} = 1, \code{ω\_down} = 1, then divides by their sum (3) because \code{congestion\_normalisation} is on:

{\fontsize{9.0}{10.8}\selectfont
\begin{verbatim}
cell     C_local  C_aisle  C_down  C_i(v)  −μ·C
-------  -------  -------  ------  ------  ------
(9,20)   0.091    0.25     0       0.114   −0.114
(10,19)  0.091    0        0.2     0.097   −0.097
(10,20)  0.077    0        0.4     0.159   −0.159
(10,21)  0.091    0        0.4     0.164   −0.164
(11,20)  0.182    0.25     0.4     0.277   −0.277
\end{verbatim}
}

Normalisation is what makes this a modulator rather than a rival. The largest \code{C\_i(v)} here is 0.277, so the entire congestion term is worth at most 0.277 — against \code{α} = 10 for a single step of progress. Turn normalisation off and \code{C\_local} reverts to a raw robot count: \code{μ·C} then reaches the scale of \code{α·Δ} (measured mean 3.40, p90 5.75, max 9.60 on this map at 40 robots), and congestion stops modulating the smaller terms and starts overruling the largest one.
\end{workedexample}
% /worked-example

\subsection{Route-level congestion}
\label{sec:congestion-route}

\begin{equation}
  \mathrm{route\_congestion}(s, g)
    = \frac{\text{occupied vertices on the route}}{|\mathrm{route}|}
    + \frac{\sum_{a \in \mathcal{A}(\mathrm{route})} \mathrm{aisle\_load}(a)}
        {\max(1, |\mathcal{A}(\mathrm{route})|)},
\end{equation}
where $\mathcal{A}(\mathrm{route})$ is the set of distinct aisles the route
touches. Used only by the assignment cost (\Cref{sec:assignment}), never by
per-step candidate scoring.
